% TOP_PROBLEM: {"id": 3108, "title": "Shuffle-Exchange Conjecture (graph-theoretic form)", "category_id": 3, "category": {"id": 3, "name": "graph_theory", "display_name": "Graph Theory", "description": "Problems involving graphs, networks, and their properties.", "slug": "graph-theory", "order_index": 3, "created_at": "2026-07-31T15:26:25.671Z"}, "status": "open"}
% FIRST_POSTED: null
% ENTRY_KIND: statement_only
% Imported from: batch01.tex; UnsolvedMath ID 3108
\documentclass[11pt]{book}
\usepackage{fontspec}
\setmainfont{Latin Modern Roman}
\setsansfont{Latin Modern Sans}
\usepackage{amsmath,amssymb,mathtools}
\usepackage{unicode-math}
\setmathfont{Latin Modern Math}
\usepackage{xurl}
\usepackage[hidelinks]{hyperref}
\newcommand{\sref}[2]{\hyperlink{source:#1:#2}{[S#2]}}
\newcommand{\cataloguescope}[1]{\paragraph{Mathematical question.}#1}
\begin{document}
% UnsolvedMath ID 3108; source catalogue code OPG-37089
\setcounter{section}{3107}
\section{Shuffle-Exchange Conjecture (graph-theoretic form)}\label{3108}
{\small\sffamily UnsolvedMath ID 3108 · Graph Theory}
\subsection{Definitions and mathematical statement}

\paragraph{Shared notation.}$\mathbb N=\{1,2,\ldots\}$, and $\mathbb Z,\mathbb Q,\mathbb R,\mathbb C$ denote the integers, rationals, reals and complex numbers. Any other conventions are specified locally.


Let $k,n\geq2$, put $t=k^{n-1}$, and let $\operatorname{SE}(k,n)$ be the bipartite graph with ordered parts $U=\{u_0,\ldots,u_{t-1}\}$ and $V=\{v_0,\ldots,v_{t-1}\}$, in which $u_iv_j$ is an edge exactly when the least nonnegative residue of $j-ki$ modulo $t$ lies in $\{0,\ldots,k-1\}$. For $r\geq2$, form an $r$-stage graph by concatenating $r-1$ identical copies, identifying output vertex $v_i$ of each copy with input vertex $u_i$ of the next. A mask is any $k$-regular bipartite multigraph between the first and last parts. The network is rearrangeable if every mask can be routed by mutually edge-disjoint paths, one for each mask edge, with the prescribed endpoints and one vertex in each consecutive stage. Define $r(k,n)$ as the smallest such $r$.
\paragraph{Statement recorded in the supplied source.}
Determine $r(k,n)$ for all $k,n\geq2$. In particular, is the proposed identity
\[r(k,n)=2n-1\]
valid for every $k,n\geq2$? \sref{3108}{2}
\subsection{Short English statement}
Determine the smallest number of identical shuffle-exchange stages that can route every prescribed k-regular end-to-end demand multigraph without sharing edges; test the 2n−1 proposal.
\subsection{Sources}
\begin{description}
\item[\hypertarget{source:3108:1}{S1}] ULAM.AI and the UnsolvedMath Contributors, \emph{UnsolvedMath: A Curated Collection of Open Mathematics Problems}, record 3108 (OPG-37089). CC BY 4.0. \url{https://huggingface.co/datasets/ulamai/UnsolvedMath}.
\item[\hypertarget{source:3108:2}{S2}] Beneš, Václav E.; Folklore; Stone, Harold S., Shuffle-Exchange Conjecture (graph-theoretic form), Open Problem Garden, node 37089. Complete problem, discussion and bibliography (including mathematical image alt text). \url{https://www.openproblemgarden.org/op/shuffle_exchange_conjecture_graph_theoretic_form}.
\end{description}
\label{end:3108}
\end{document}
